\documentclass[11pt,a4paper]{article}
\usepackage[T1]{fontenc}
\usepackage{amsmath,amssymb,geometry}
\geometry{margin=28mm}
\title{EPPA--III: a short orbit proof for Lemma 5.12}
\author{Independent validation supplement; manuscript unchanged}
\date{10 October 2026}
\begin{document}
\maketitle

\paragraph{Source anchor.}
This note addresses the final orbit-transitivity assertion
of Lemma~\texttt{lemma:doublecover}, immediately after
Claim~\texttt{claim:taylorconclusion} in
\texttt{main-new(8).tex} (new Section~5.3).
The manuscript also proves a more elaborate claim
\texttt{claim:orbit} using a maximal partial
automorphism. The shorter route below is independent
of that maximal-domain argument.

\paragraph{Assumptions.}
Let \(G\) be a finite non-homogeneous graph and
\(H\) a transitive imprimitive EPPA witness with
\(|H|=2|G|\).
Assume the invariant two-element block system
\(P=\{\{x_0,x_1\}:x\in V(G)\}\) has already been
established. By the separate four-Boolean
matrix argument, the inter-block adjacency is
that of the Taylor double, with a matching
between mates uniformly present or absent.

\paragraph{Short proof of the quotient action.}
Let \(K\le \operatorname{Sym}(V(G))\) be the
permutation group induced by \(\operatorname{Aut}(H)\)
on the invariant blocks. By the two-point EPPA
property, \(K\) is transitive on ordered pairs
\((x,y)\) with \(x\sim_G y\), and separately on
ordered pairs with \(x\not\sim_G y\), \(x\ne y\).
If these were distinct \(K\)-orbits, every
element of \(K\) would preserve adjacency in
\(G\); indeed \(K\le \operatorname{Aut}(G)\).
Projecting any EPPA extension onto the blocks
would then extend the original partial
automorphism within \(G\). Consequently \(G\)
would be homogeneous, a contradiction.
Thus the two orbits coincide, and \(K\)
is 2-transitive.

\paragraph{Lifting to ordered edges and nonedges.}
Between any two different blocks \(\{x_0,x_1\}\)
and \(\{y_0,y_1\}\), the Taylor pattern has
exactly two edges and two nonedges.
The deck involution simultaneously swapping
\(x_0\) with \(x_1\) for every \(x\) is an
automorphism, even if the matching inside
each block is present. This involution
exchanges the two ordered edges over a fixed
ordered pair of blocks, and likewise their
two nonedges.
Given any two ordered inter-block edges,
2-transitivity of \(K\) first aligns their
ordered blocks, and the deck involution
then aligns the two ordered vertex pairs.
The same argument works for nonedges.
This is precisely the remaining assertion
of Lemma~3.3.

\paragraph{Formal verification.}
The quotient 2-transitivity argument is
verified in \texttt{DoubleCover/FibreOrbits.lean};
the deck involution and pair-status
classification in \texttt{DoubleCover/DeckSwap.lean};
and the canonical ordered inter-fibre
transitivity in \texttt{DoubleCover/InterblockOrbits.lean}.
\texttt{BlockSystems/OrderedInterblockOrbits.lean}
transports the orbit conclusion to the
original host, with explicit invariance,
EPPA and size hypotheses.
The complete statements are axiom-audited.
The use of Gardiner's classification to
deduce non-homogeneity from the exact list
of excluded graphs in the source lemma is
still a paper-level dependency and is not
implicitly proved by this Lean development.

\paragraph{Editorial status.}
This is an optional proof simplification
and validation supplement, not an instruction
to rewrite the existing manuscript.
Retain the proposed Claim~3.6 clarification
allowing uniform mate edges.
\end{document}
